% !PRINTED 338
\weberpage{338}
% running head: 338 | Elfter Abschnitt. | §. 91.

Um $\varDelta_{\lambda,\mu}$ für die canonische Form zu berechnen, hat man
in \S.\,90, (3):
\[
h, \qquad m
\]
durch
\[
h(\lambda-6m^2\mu), \qquad m\lambda+(h^3+2m^3)\mu
\]
zu ersetzen, und man findet so:
\[
\varDelta_{\lambda,\mu} = -\,h(\lambda-6m^2\mu)\,[m\lambda+(h^3+2m^3)\mu]^2\,\Sigma y_1^3
+ \{h^3(\lambda-6m^2\mu)^3+2[m\lambda+(h^3+2m^3)\mu]^3\}\,y_1y_2y_3,
\]
woraus nach \S.\,90, (7), (8) mit Rücksicht auf (2) nach einigen
einfachen Rechnungen folgt:
\weq{5}{\begin{gathered}
L = -2\lambda^2\mu(h^3-m^3)m - \lambda\mu^2(h^6-20h^3m^3-8m^6) \\
\qquad{}+ 8\mu^3m^2(h^3-m^3)^2, \\
M = \lambda^3+12\lambda\mu^2m(h^3-m^3)+2\mu^3(h^6-20h^3m^3-8m^6).
\end{gathered}}

Man bekommt also auf diese Weise nur zwei Invarianten vom
$4^{\mathrm{ten}}$ und $6^{\mathrm{ten}}$ Grade, nämlich in der canonischen Form:
\weq{6}{m(h^3-m^3), \quad h^6-20h^3m^3-8m^6.}

Das Gewicht dieser Invarianten ist [nach Bd. I, \S.\,60, (3)]
gleichfalls 4 und 6, und wenn wir sie also in der allgemeinen
Form mit $S$ und $T$ bezeichnen und $h=1$ setzen, so findet sich
\weq{7}{\begin{gathered}
m(1-m^3) = r^4S, \\
1-20m^3-8m^6 = r^6T.
\end{gathered}}

Für $L$ und $M$ erhält man aus (5) die Darstellung:
\weq{8}{\begin{gathered}
L = -2S\lambda^2\mu - T\lambda\mu^2 + 8S^2\mu^3 \\
M = \lambda^3+12S\lambda\mu^2+2T\mu^3.
\end{gathered}}

Die Functionen $S$ und $T$ vom $4^{\mathrm{ten}}$ und $6^{\mathrm{ten}}$ Grade sind
zuerst von \gesperrt{Aronhold} berechnet worden. Wir wollen die langen
Ausdrücke nicht hierher setzen, betrachten aber diese beiden
Grössen jetzt als bekannt und gegeben.

Aus (7) lässt sich eine biquadratische Gleichung ableiten
für das Verhältniss $m^2:r^2$. Man findet nämlich daraus:
\weq{9}{\begin{gathered}
m^5-m^8 = r^4m^4S \\
m^2-2m^5+m^8 = r^8S^2 \\
m^2-20m^5-8m^8 = r^6m^2T.
\end{gathered}}

Wenn man daraus $m^2$ und $m^5$ eliminirt, so findet man
\weq{10}{27m^8+18Sm^4r^4+Tm^2r^6-S^2r^8=0.}
